\chapter{Resultados Esperados e Cronograma de Atividades}
\label{cap:cronograma}

\section{Resultados esperados}

O modelo descrito no Capítulo~\ref{cap:modelagem} encontra-se verificado e
operacional, e os resultados do Capítulo~\ref{cap:resultados} estabelecem o
regime de operação do sistema. As etapas restantes do trabalho destinam-se a
substituir as estimativas ainda empregadas por dados medidos, a explorar o
espaço de parâmetros de projeto e a confrontar o modelo com a resposta
experimental do reator.

\subsection{Substituição dos parâmetros estimados}

Três grandezas do modelo são hoje estimativas, e cada uma delas será substituída
por medida experimental assim que disponível.

A massa específica e a espessura do \textit{washcoat} entram no modelo apenas
pelo produto $\rho_{wc}\,\delta_{wc}$, ao qual a taxa de reação é linearmente
proporcional, conforme as Equações~\eqref{eq:rw} e~\eqref{eq:w}. Sua
caracterização é, por essa razão, a medida de maior retorno imediato para a
acurácia do modelo, e espera-se que a incerteza associada ao número de Damköhler
caia proporcionalmente.

As dimensões do canal, hoje tomadas de valores típicos de monólitos comerciais,
serão medidas nas peças efetivamente disponíveis no laboratório. A altura e o
comprimento do canal deixaram de ser parâmetros secundários a partir do momento
em que se verificou que o comprimento de entrada mássica supera o comprimento
físico em quase duas ordens de grandeza; espera-se que a razão $L_e/L$ se
mantenha elevada, mas seu valor exato determina a magnitude da limitação
geométrica discutida no Capítulo~\ref{cap:resultados}.

As propriedades físicas do meio reacional, por fim, serão obtidas para a carga
efetivamente empregada, em lugar dos valores de literatura hoje adotados.

\subsection{A geometria do monólito disponível}
\label{sec:geomreal}

Durante a redação deste texto, as dimensões das peças de suporte efetivamente
disponíveis no laboratório foram medidas e repassadas pela orientação. Diferem
de forma apreciável da geometria de referência adotada nas simulações
apresentadas, e merecem registro por condicionarem as etapas seguintes.

\begin{table}[htb]
\centering
\caption{Geometria simulada e geometria das peças disponíveis}
\label{tab:geomreal}
\begin{tabular}{l c c}
\toprule
 & Simulado & Peça disponível \\
\midrule
Seção do canal          & plano          & quadrado \\
Abertura                & \SI{1.1}{\milli\meter} & \SI{0.55}{\milli\meter} \\
Comprimento             & \SI{50}{\milli\meter}  & \SI{29.15}{\milli\meter} \\
Paredes catalíticas     & \num{1}        & \num{4} \\
Espessura de parede     & ---            & \SI{0.10}{\milli\meter} \\
Densidade de células    & \num{400} cpsi & \num{1527} cpsi \\
\bottomrule
\end{tabular}
\end{table}

A geometria simulada derivou de \textcite{Asli2011}, como exposto no
Capítulo~\ref{cap:revisao}. A peça disponível tem metade da abertura, seção
quadrada com as quatro paredes recobertas e densidade de células quase quatro
vezes maior.

Estimativas preliminares, obtidas por extrapolação do modelo verificado através
da solução de Lévêque e do balanço de área, indicam que o canal menor engrossa
menos a camada-limite e eleva o coeficiente de transferência de massa em cerca
de \SI{50}{\percent}. O número de Damköhler cairia de \num{0.76} para
aproximadamente \num{0.51}, permanecendo em regime misto, e a conversão por
passe de uma peça subiria de \SI{1.4}{\percent} para algo próximo de
\SI{7.7}{\percent}, sobretudo pelo aumento de oito vezes na razão entre área de
parede e volume. O comprimento de entrada mássica cairia para
\SI{216}{\milli\meter}, ainda sete vezes o comprimento da peça.

Dessa última razão decorre uma consequência de projeto. Operando na região de
entrada, o coeficiente de transferência de massa decai com $x^{-1/3}$, de modo
que empilhar peças curtas com desalinhamento entre elas, de forma que a
camada-limite se reinicie a cada peça, eleva o coeficiente médio por um fator
$N^{1/3}$ em relação a um canal contínuo de mesmo comprimento total. Para trinta
peças, fator próximo de três. A verificação dessa estimativa por simulação da
geometria real consta das atividades previstas.

\subsection{Varredura paramétrica}

Com o modelo calibrado, pretende-se realizar varreduras em três direções.

A espessura do \textit{washcoat} será variada de modo a localizar o joelho da
curva de saturação descrita pela Equação~\eqref{eq:saturacao}, que define a
faixa de deposição economicamente razoável. A análise preliminar da
Seção~\ref{sec:thiele} indica \SIrange{10}{15}{\micro\meter}, e espera-se que
a varredura confirme essa faixa com o valor medido de $D_p$ e forneça um número
que possa orientar o procedimento de impregnação.

A temperatura será varrida na faixa compatível com a operação em metanol à
pressão atmosférica. Essa varredura exige a substituição da constante cinética
por uma expressão de Arrhenius, o que depende da energia de ativação a ser
obtida experimentalmente. Espera-se observar o deslocamento do número de
Damköhler com a temperatura, dado que a cinética é mais sensível a $T$ do que a
difusividade, e portanto a transição gradual do regime misto em direção ao
controle por transporte à medida que a temperatura se eleva.

A geometria do canal será reconstruída segundo as dimensões da
Tabela~\ref{tab:geomreal}, o que exige abandonar a formulação pseudo-bidimensional
em favor de um quarto de duto quadrado, aproveitando os dois planos de simetria.
A física, as funções de campo e o protocolo de verificação do
Capítulo~\ref{cap:modelagem} transferem-se sem alteração; mudam a geometria e a
malha. Em seguida, comprimento e abertura serão variados em torno desses valores,
e o efeito do desalinhamento entre peças sucessivas será avaliado.

\subsection{Calibração e confronto com dados experimentais}

Os ensaios em condução no laboratório fornecerão conversões medidas em condições
controladas de temperatura, razão molar e vazão. O confronto entre essas
conversões e as previstas pelo modelo permitirá ajustar os parâmetros ainda
incertos, em particular a difusividade efetiva do triglicerídeo no meio e no
poro, e estabelecer a faixa de validade do modelo.

Espera-se que o modelo reproduza a tendência experimental de conversão com a
vazão, dado que essa dependência é governada pelo transporte externo, que o
modelo resolve explicitamente. Discrepâncias sistemáticas, caso ocorram, serão
examinadas à luz das hipóteses listadas na Seção~\ref{sec:difusividades},
notadamente a de difusividades iguais entre espécies, cuja substituição pela
formulação de Maxwell-Stefan constitui a alternativa prevista.

\subsection{Decisão quanto à rota alcoólica}

Permanece em aberto a definição da rota a ser modelada em definitivo. A rota
metílica, à qual correspondem os parâmetros de \textcite{Melo2023} empregados
até aqui, opera a \SI{65}{\celsius} com razão molar de \num{4}:\num{1}. A rota
etílica, estudada por \textcite{Fonseca2026} a \SI{70}{\celsius} e razão de
\num{1}:\num{9}, apresenta vantagens de origem da matéria-prima no contexto
brasileiro, mas cinética mais lenta e separação mais difícil. A decisão será
tomada em conjunto com a orientação, e o modelo está estruturado de modo a
acomodar qualquer das duas sem alteração de sua formulação.

\section{Cronograma de atividades}

A Tabela~\ref{tab:cronograma} apresenta o cronograma proposto para a execução do
trabalho, considerando-se o período integral do mestrado.

\begin{table}[htb]
\centering
\caption{Cronograma proposto para as atividades do presente estudo}
\label{tab:cronograma}
\renewcommand{\arraystretch}{1.3}
\begin{tabular}{|l|c|c|c|c|c|c|}
\hline
\multirow{2}{*}{\textbf{Atividade}} & \multicolumn{6}{c|}{\textbf{Semestre / Ano}} \\
\cline{2-7}
 & \textbf{1º/25} & \textbf{2º/25} & \textbf{1º/26} & \textbf{2º/26} & \textbf{1º/27} & \textbf{2º/27} \\
\hline
Disciplinas obrigatórias        & \cellcolor{gray!55} & \cellcolor{gray!55} & \cellcolor{gray!55} & \cellcolor{gray!55} &  &  \\
\hline
Revisão bibliográfica           & \cellcolor{gray!55} & \cellcolor{gray!55} & \cellcolor{gray!55} & \cellcolor{gray!55} & \cellcolor{gray!55} &  \\
\hline
Desenvolvimento do modelo CFD   &  & \cellcolor{gray!55} & \cellcolor{gray!55} & \cellcolor{gray!55} &  &  \\
\hline
Verificação do modelo           &  &  & \cellcolor{gray!55} & \cellcolor{gray!55} &  &  \\
\hline
Exame de qualificação           &  &  &  & \cellcolor{gray!55} &  &  \\
\hline
Caracterização do catalisador   &  &  &  & \cellcolor{gray!55} & \cellcolor{gray!55} &  \\
\hline
Ensaios experimentais           &  &  &  & \cellcolor{gray!55} & \cellcolor{gray!55} &  \\
\hline
Calibração do modelo            &  &  &  &  & \cellcolor{gray!55} &  \\
\hline
Varredura paramétrica           &  &  &  &  & \cellcolor{gray!55} & \cellcolor{gray!55} \\
\hline
Análise dos resultados          &  &  &  &  & \cellcolor{gray!55} & \cellcolor{gray!55} \\
\hline
Redação e defesa da dissertação &  &  &  &  &  & \cellcolor{gray!55} \\
\hline
\end{tabular}
\end{table}

\pendente{conferir as datas de início do mestrado e de integralização dos
créditos, e ajustar o cronograma à previsão de defesa acordada com a orientação}
